\documentclass[12pt,a4paper]{article}
\usepackage[utf8]{inputenc}
\usepackage[bulgarian]{babel}
\usepackage{amsmath}
\usepackage{amsfonts}
\usepackage{amssymb}
\usepackage{graphicx}
\usepackage{geometry}
\usepackage{fancyhdr}
\usepackage{xcolor}

\geometry{margin=2cm}
\pagestyle{fancy}
\fancyhead[L]{Проект MORT: Техническа Спецификация}
\fancyhead[R]{Версия 2.1}

\title{\textbf{Пълна Архитектурна и Теоретична Обосновка на MORT} \\ \large (Micro-Optotronic Reactive Tesselator) \\ \small Хибридна самоорганизираща се система за вълнова теселация}
\author{Отдел по Авангардна Инженерна Логика}
\date{2026 г.}

\begin{document}

\maketitle

\begin{abstract}
Настоящият документ представлява детайлно описание на компютърната архитектура MORT. Системата заменя традиционния транзисторен превключвател с принципа на „Прожектираната регенерация“. Изчисленията се извършват чрез взаимодействието на терахерцови електромагнитни вълни с динамично конфигурируеми топологични полета. Докладът разглежда решенията за синхронизация, термична устойчивост, мрежова комуникация и енергийна рекуперация.
\end{abstract}

\tableofcontents
\newpage

\section{Фундаментална Концепция: Теселация срещу Изчисление}
Традиционните компютри обработват данни чрез последователно превключване на логически порти. MORT използва \textbf{пространствена интерференция}. Изчислението е резултат от преминаването на сложен вълнов фронт през специфична геометрия (Shield), при което се получава теселирано решение в реално време.

\section{Изчислително Ядро и Динамична Топология}

\subsection{Електромагнитният Щит (The Yin-Yang Shield)}
Вместо статичен метален филтър, централното ядро е \textbf{електромагнитна структура във време-пространството}. 
\begin{itemize}
    \item \textbf{Принцип:} Полетата се модулират така, че да създават зони с различна плътност и фазово съпротивление.
    \item \textbf{Предимство:} Липса на физическа „умора“ на материала. Структурата се „плете“ от контролирани стоящи вълни, което позволява моментално преконфигуриране на архитектурата (от аритметично към логическо ядро).
\end{itemize}

\subsection{D-куплунг (Detector-Coupler): Регенеративната Връзка}
За да се преодолее ентропията, MORT използва активни \textbf{D-куплунги}. Всеки куплунг функционира като фото-електронен репитер:
\begin{equation}
    \Psi_{out} = G \cdot \Psi_{in} e^{i\phi_{offset}}
\end{equation}
Където $G$ е коефициентът на усилване, а $\phi_{offset}$ е програмираното фазово отместване. Това гарантира, че сигналът остава кохерентен без значение от броя преминали ядра.

\section{Синхронизация и Системна Стабилност}

\subsection{M.A.M. (Monophotonic Atomic Marker)}
Вълновият синхрон изисква точност под $10^{-12}$ секунди. М.А.М. модулът използва един фотон, затворен в идеален оптичен контур. Всеки цикъл на фотона генерира импулс, който служи за „сърцебиене“ на системата. Това е \textbf{квантов метроном}, който гарантира, че „песента“ (данните) и „диригентът“ (часовникът) са винаги в синхрон.

\subsection{RF Backflip Detector: Имунната Система}
Всяка вълнова система страда от паразитни отражения (ехо). \textbf{Backflip Detector} е многочестотен сензор, който следи за „обратни вълни“. 
\begin{itemize}
    \item \textbf{Функция:} При засичане на аномалия, детекторът подава \textbf{изправителен сигнал}. Това е реактивен офсет, който неутрализира шума чрез деструктивна интерференция, преди той да корумпира данните.
\end{itemize}

\subsection{Мексиканска Вълна (System Refresh)}
За разлика от грубия „рестарт“, MORT използва \textbf{Мексиканска вълна}. Това е последователен вълнов фронт, който преминава през всички модули (ПАМ, Ядра, Трансрутери). Той „продухва“ натрупаните фазови остатъци и калибрира системата в движение. Скоростта на MORT позволява това да се случва хиляди пъти в секунда, без потребителят да усети прекъсване.

\section{Мрежова Архитектура и ТрансБука}

\subsection{Хардуерни Диалекти}
Поради безкрайните конфигурации на полевия щит, всяка машина развива свой уникален „диалект“. Това е най-високото ниво на хардуерна защита – сигнал от една машина е неизползваем за друга без правилния ключ.

\subsection{ТрансБука (Универсален Транслатор)}
За комуникация се използват два типа транслатори:
\begin{itemize}
    \item \textbf{Изходен:} Конвертира локалния диалект в „Златен стандарт“.
    \item \textbf{Входен:} Прилага специфичния за приемащата машина офсет.
    \item \textbf{Пилотна вълна:} Пътува заедно с данните като еталон. Ако средата (кабела) изкриви вълната, пилотната вълна показва степента на отклонение, което позволява на AI агента на приемника да коригира сигнала в реално време.
\end{itemize}

\section{Интеграция на AI Агенти (Native Data Interpreters)}
AI не е софтуер над MORT – той е \textbf{неразделна част от интерфейса}. 
\begin{itemize}
    \item \textbf{Роля:} AI агентите действат като преводачи на вълновата симфония. Те се самообучават върху физическите особености (дефекти и отклонения) на конкретното ядро и се научават да ги използват за оптимизация на изчисленията.
    \item \textbf{Защита:} AI разпознава „органичния ритъм“ на входящите данни. Всеки опит за външна намеса се засича като „дисхармония“ и се блокира нативно.
\end{itemize}

\section{Енергиен Баланс и Wolfram Изчисления}
MORT е проектиран като \textbf{енергийно затворена система}. 
\begin{enumerate}
    \item \textbf{Recycling:} Използваните вълни след теселация не се губят, а се конвертират обратно в захранващ импулс чрез свръхпроводникови актуатори.
    \item \textbf{Wolfram Logic:} Алгоритмите за самоизчисляване (Self-computing) позволяват на програмата хем да изпълнява логика, хем да поддържа собствения си енергиен баланс. Колкото повече задачи изпълнява системата, толкова повече кохерентна енергия се генерира чрез подреждане на информационното поле.
\end{enumerate}

\section{Заключение: МОРТ като Жив Хардуер}
Комбинацията от терахерцова физика, квантова синхронизация (M.A.M.) и AI интерпретация превръща MORT в „Жив хардуер“. Системата притежава метаболизъм (рециклиране на енергия), рефлекси (Backflip Detector) и разум (Native AI). MORT не просто пресмята бъдещето – той го плете чрез вълновата структура на реалността.

\end{document}
